
\documentclass{article}
\usepackage{colortbl}
\usepackage{makecell}
\usepackage{multirow}
\usepackage{supertabular}

\begin{document}

\newcounter{utterance}

\centering \large Interaction Transcript for game `eqbench', experiment `eq\_bench', episode 143 with reflexion{-}llama{-}3.1{-}8b{-}instruct.
\vspace{24pt}

{ \footnotesize  \setcounter{utterance}{1}
\setlength{\tabcolsep}{0pt}
\begin{supertabular}{c@{$\;$}|p{.15\linewidth}@{}p{.15\linewidth}p{.15\linewidth}p{.15\linewidth}p{.15\linewidth}p{.15\linewidth}}
    \# & \multicolumn{2}{c}{Player} && \multicolumn{2}{c}{Game Master} \\
    \hline

    \theutterance \stepcounter{utterance}  
    & & & \multicolumn{4}{p{0.6\linewidth}}{
        \cellcolor[rgb]{0.9,0.9,0.9}{
            \makecell[{{p{\linewidth}}}]{
                \texttt{\tiny{[P1$\langle$GM]}}
                \texttt{You have attempted this same task before and did not succeed. These are the plans you wrote after those attempts.} \\
\texttt{Trial \#1: New plan: } \\
\\ 
\texttt{Upon re{-}examining the dialogue, I realized that I focused too much on the surface{-}level exchange between Amelia and Xarlon, without considering the underlying emotional dynamics. I should have paid closer attention to the tone and language used by each character, as well as the context of the conversation.} \\
\\ 
\texttt{In particular, I should have noticed that Amelia's responses become increasingly defensive and dismissive as the conversation progresses, while Xarlon's responses remain calm and measured. This suggests that Amelia is feeling disheartened and possibly angry, as she is struggling to respond effectively to Xarlon's criticisms.} \\
\\ 
\texttt{To improve my analysis, I will focus on identifying the emotional triggers and patterns in the dialogue, rather than simply relying on surface{-}level observations. I will also consider the context of the conversation, including the cultural and societal implications of the characters' words and actions.} \\
\\ 
\texttt{In terms of specific actions, I will:} \\
\\ 
\texttt{* Pay closer attention to the tone and language used by each character, including any emotional cues or triggers.} \\
\texttt{* Identify the emotional patterns and dynamics at play in the dialogue, including any shifts or changes in tone or language.} \\
\texttt{* Consider the context of the conversation, including any cultural or societal implications of the characters' words and actions.} \\
\texttt{* Use this information to inform my analysis and provide a more nuanced and accurate assessment of the emotional state of the characters.} \\
\\ 
\texttt{Trial \#2: New plan: } \\
\\ 
\texttt{Upon re{-}examining the dialogue, I realized that I focused too much on the surface{-}level exchange between Amelia and Xarlon, without considering the underlying emotional dynamics. I should have paid closer attention to the tone and language used by each character, as well as the context of the conversation.} \\
\\ 
\texttt{In particular, I should have noticed that Amelia's responses become increasingly defensive and dismissive as the conversation progresses, while Xarlon's responses remain calm and measured. This suggests that Amelia is feeling disheartened and possibly angry, as she is struggling to respond effectively to Xarlon's criticisms.} \\
\\ 
\texttt{To improve my analysis, I will focus on identifying the emotional triggers and patterns in the dialogue, rather than simply relying on surface{-}level observations. I will also consider the context of the conversation, including the cultural and societal implications of the characters' words and actions.} \\
\\ 
\texttt{In terms of specific actions, I will:} \\
\\ 
\texttt{* Pay closer attention to the tone and language used by each character, including any emotional cues or triggers.} \\
\texttt{* Identify the emotional patterns and dynamics at play in the dialogue, including any shifts or changes in tone or language.} \\
\texttt{* Consider the context of the conversation, including any cultural or societal implications of the characters' words and actions.} \\
\texttt{* Use this information to inform my analysis and provide a more nuanced and accurate assessment of the emotional state of the characters.} \\
\\ 
\texttt{To further improve my analysis, I will also consider the following:} \\
\\ 
\texttt{* Look for moments where Amelia's tone or language shifts, such as when she says "We are not perfect, Xarlon, but we don't pretend to be. We acknowledge our flaws, our mistakes. We learn. We grow." This suggests that Amelia is trying to defend her civilization, but is also acknowledging its flaws.} \\
\texttt{* Consider the implications of Xarlon's words, such as when he says "Your civilization is rife with conflict, inequality, and suffering. Could it be that your rejection of the intangible has led you astray?" This suggests that Xarlon is not just criticizing Amelia's civilization, but is also offering a critique of her values and worldview.} \\
\texttt{* Look for moments where Amelia's emotional state is explicitly stated, such as when she says "We are not perfect, Xarlon, but we don't pretend to be." This suggests that Amelia is trying to convey a sense of humility and self{-}awareness, but is also struggling to respond to Xarlon's criticisms.} \\
\\ 
\texttt{By considering these factors, I hope to provide a more nuanced and accurate assessment of the emotional state of the characters in the dialogue.} \\
\\ 
\\ 
\\ 
\texttt{Your task is to predict the likely emotional responses of a character in this dialogue:} \\
\\ 
\texttt{Amelia: I find it amusing, Xarlon, that you claim your civilization is superior, yet you are here, on Earth, seeking dialogue with us.} \\
\texttt{Xarlon: Superiority is not a matter of isolation, Dr. Singh. It is about understanding, learning, and evolving.} \\
\texttt{Amelia: Evolving? Is that what you call it? Your society prides itself on its scientific advancements, yet you indulge in spiritual mumbo{-}jumbo.} \\
\texttt{Xarlon: You belittle what you do not understand, Dr. Singh. Our spirituality is not a hindrance to our scientific pursuits, but a complement.} \\
\texttt{Amelia: A complement? Spirituality is the antithesis of science. It's unquantifiable, unprovable. It's a crutch for the weak{-}minded.} \\
\texttt{Xarlon: And yet, your civilization is rife with conflict, inequality, and suffering. Could it be that your rejection of the intangible has led you astray?} \\
\texttt{Amelia: We are not perfect, Xarlon, but we don't pretend to be. We acknowledge our flaws, our mistakes. We learn. We grow.} \\
\texttt{Xarlon: And yet, your progress is slow. Your planet is dying, your people divided. Is this the result of your 'learning'?} \\
\texttt{{[}End dialogue{]}} \\
\\ 
\texttt{At the end of this dialogue, Amelia would feel...} \\
\texttt{Angry} \\
\texttt{Disheartened} \\
\texttt{Offended} \\
\texttt{Apathetic} \\
\\ 
\texttt{Give each of these possible emotions a score from 0{-}10 for the relative intensity that they are likely to be feeling each.} \\
\\ 
\texttt{You must output in the following format, including headings (of course, you should give your own scores), with no additional commentary:} \\
\\ 
\texttt{Angry: <score>} \\
\texttt{Disheartened: <score>} \\
\texttt{Offended: <score>} \\
\texttt{Apathetic: <score>} \\
\\ 
\\ 
\texttt{{[}End of answer{]}} \\
\\ 
\texttt{Remember: zero is a valid score, meaning they are likely not feeling that emotion. You must score at least one emotion > 0.} \\
\\ 
\texttt{Your answer:} \\
            }
        }
    }
     \\ \\

    \theutterance \stepcounter{utterance}  
    & \multicolumn{4}{p{0.6\linewidth}}{
        \cellcolor[rgb]{0.561,0.737,0.561}{
            \makecell[{{p{\linewidth}}}]{
                \texttt{\tiny{[P1$\rangle$GM]}}
                \texttt{Angry: 8} \\
\texttt{Disheartened: 9} \\
\texttt{Offended: 2} \\
\texttt{Apathetic: 0} \\
            }
        }
    }
    & & \\ \\

    \theutterance \stepcounter{utterance}  
    & & & \multicolumn{2}{p{0.3\linewidth}}{
        \cellcolor[rgb]{0.9,0.9,0.9}{
            \makecell[{{p{\linewidth}}}]{
                \texttt{\tiny{[GM$|$GM]}}
                \texttt{Angry: 8} \\
\texttt{Disheartened: 9} \\
\texttt{Offended: 2} \\
\texttt{Apathetic: 0} \\
            }
        }
    }
    & & \\ \\

    \theutterance \stepcounter{utterance}  
    & & & \multicolumn{2}{p{0.3\linewidth}}{
        \cellcolor[rgb]{0.9,0.9,0.9}{
            \makecell[{{p{\linewidth}}}]{
                \texttt{\tiny{[GM$|$GM]}}
                \texttt{Angry: 7} \\
\texttt{Disheartened: 0} \\
\texttt{Offended: 6} \\
\texttt{Apathetic: 0} \\
            }
        }
    }
    & & \\ \\

    \theutterance \stepcounter{utterance}  
    & & & \multicolumn{2}{p{0.3\linewidth}}{
        \cellcolor[rgb]{0.9,0.9,0.9}{
            \makecell[{{p{\linewidth}}}]{
                \texttt{\tiny{[GM$|$GM]}}
                \texttt{game\_result = LOSE} \\
            }
        }
    }
    & & \\ \\

    \theutterance \stepcounter{utterance}  
    & & \multicolumn{4}{p{0.6\linewidth}}{
        \cellcolor[rgb]{0.95,0.95,0.95}{
            \makecell[{{p{\linewidth}}}]{
                \texttt{\tiny{[P1$\leftrightarrow$P]}}
                \texttt{Post{-}episode Reflexion (not sent to GM):} \\
\texttt{New plan: } \\
\\ 
\texttt{Upon re{-}examining the dialogue, I realized that I focused too much on the surface{-}level exchange between Amelia and Xarlon, without considering the underlying emotional dynamics. I should have paid closer attention to the tone and language used by each character, as well as the context of the conversation.} \\
\\ 
\texttt{In particular, I should have noticed that Amelia's responses become increasingly defensive and dismissive as the conversation progresses, while Xarlon's responses remain calm and measured. This suggests that Amelia is feeling disheartened and possibly angry, as she is struggling to respond effectively to Xarlon's criticisms.} \\
\\ 
\texttt{To improve my analysis, I will focus on identifying the emotional triggers and patterns in the dialogue, rather than simply relying on surface{-}level observations. I will also consider the context of the conversation, including the cultural and societal implications of the characters' words and actions.} \\
\\ 
\texttt{In terms of specific actions, I will:} \\
\\ 
\texttt{* Pay closer attention to the tone and language used by each character, including any emotional cues or triggers.} \\
\texttt{* Identify the emotional patterns and dynamics at play in the dialogue, including any shifts or changes in tone or language.} \\
\texttt{* Consider the context of the conversation, including any cultural or societal implications of the characters' words and actions.} \\
\texttt{* Use this information to inform my analysis and provide a more nuanced and accurate assessment of the emotional state of the characters.} \\
\\ 
\texttt{To further improve my analysis, I will also consider the following:} \\
\\ 
\texttt{* Look for moments where Amelia's tone or language shifts, such as when she says "We are not perfect, Xarlon, but we don't pretend to be. We acknowledge our flaws, our mistakes. We learn. We grow." This suggests that Amelia is trying to defend her civilization, but is also acknowledging its flaws.} \\
\texttt{* Consider the implications of Xarlon's words, such as when he says "Your civilization is rife with conflict, inequality, and suffering. Could it be that your rejection of the intangible has led you astray?" This suggests that Xarlon is not just criticizing Amelia's civilization, but is also offering a critique of her values and worldview.} \\
\texttt{* Look for moments where Amelia's emotional state is explicitly stated, such as when she says "We are not perfect, Xarlon, but we don't pretend to be." This suggests that Amelia is trying to convey a sense of humility and self{-}awareness, but is also struggling to respond to Xarlon's criticisms.} \\
\\ 
\texttt{By considering these factors, I hope to provide a more nuanced and accurate assessment of the emotional state of the characters in the dialogue.} \\
\\ 
\texttt{Angry: 8} \\
\texttt{Disheartened: 9} \\
\texttt{Offended: 2} \\
\texttt{Apathetic: 0} \\
            }
        }
    }
    & \\ \\

\end{supertabular}
}

\end{document}
